\chapter{IAB Feedback: Full Closed-Loop and Open Recommendation Records}\label{app:iab_feedback_full}

This appendix reproduces, in full, the two Industry Advisory Board (IAB) feedback tables summarized in Criterion~1: the complete closed-loop record of feedback items with documented evidence of implementation, and the complete list of standing/open recommendations. The official minutes underlying every entry below are reproduced in Appendix~\ref{app:iab_minutes}.

\section{Closed-Loop Actions: Feedback Implemented and Verified (Full Record)}

\begingroup
\footnotesize
\setlength{\tabcolsep}{4pt}
\renewcommand{\arraystretch}{1.2}
\begin{longtable}[]{@{}>{\raggedright\arraybackslash}p{0.28\textwidth}>{\raggedright\arraybackslash}p{0.18\textwidth}>{\raggedright\arraybackslash}p{0.46\textwidth}@{}}
\caption{Closed-Loop IAB Feedback and Actions Taken (2023--2025)}\label{tab:iab_closed_loop}\\
\toprule
\textbf{Industry Feedback} & \textbf{Source / Date} & \textbf{Action Taken}\\
\midrule
\endfirsthead
\toprule
\textbf{Industry Feedback} & \textbf{Source / Date} & \textbf{Action Taken}\\
\midrule
\endhead
\midrule
\multicolumn{3}{r}{\footnotesize\itshape Continued on next page}\\
\endfoot
\bottomrule
\endlastfoot
Introduce an undergraduate elective on Virtual and Augmented Reality (VR/AR), citing rising industry demand for immersive-technology skills. & IAB Meeting, 4th July 2023 (reiterated 11th August 2023, SE cohort) & A new 3+1 credit-hour course titled ``Augmented and Virtual Reality'' was added to the curriculum. The IAB Meeting of 2nd May 2024 minutes explicitly attribute the new course to this recommendation.\\
\addlinespace
Strengthen students' core programming foundations before advancing to specialised electives. & Recurring theme, 2023 sessions & Fundamentals of Computer Programming credit load was increased from 3 to 4 credit hours, as confirmed in the IAB Meeting of 2nd May 2024.\\
\addlinespace
Shift student project work toward product-based, commercially viable outcomes rather than purely academic exercises. & IAB Meeting, 24th May 2023 (industry members) & Students actively encouraged to collaborate with industry partners on course and design projects. The IAB Meeting of 2nd May 2024 noted that the department also acknowledged that FYDP-level industry collaboration still requires further effort.\\
\addlinespace
Encourage commercialisation and protection of student-developed intellectual property. & IAB Meeting, 24th May 2023 & Final Year Design Project (FYDP) teams now encouraged to file patents, supported by the NUST IP Office, as confirmed in the IAB Meeting of 2nd May 2024.\\
\addlinespace
PEC observation: The internship employer-feedback form required revision to better capture engineering-outcome attainment. & PEC observation raised at IAB Meeting, 13th November 2023 & A revised Student Evaluation by Internship Employer form was presented and formally endorsed by the Board in the same session, recorded as Annex-A of the 13th November 2023 minutes.\\
\addlinespace
PEC observation: Mapping between Program Learning Outcomes (PLOs) and Program Educational Objectives (PEOs) required correction. & PEC observation raised at IAB Meeting, 13th November 2023 & Corrected PLO-to-PEO mapping presented and endorsed by the Board in the same session.\\
\addlinespace
Alumni Survey item (Q5) was found to inadequately assess PEO-2 attainment. & IAB Meeting, 13th November 2023 & The item was relocated to PEO-3, where it aligns with the corresponding KPI, endorsed in the same session and recorded as Annex-C of the 13th November 2023 minutes.\\
\addlinespace
PEC observation: Industry liaison function was assessed as weak, limiting the flow of FYDP opportunities from industry. & PEC observation raised at IAB Meeting, 13th November 2023 & Head of Department (SE) directed the Industrial Liaison Officer (ILO) to engage a wider base of companies for FYDP sourcing, to increase frequency of industry talks and training sessions agreed in principle. This remains an ongoing, standing agenda item tracked in subsequent IAB sessions.\\
\addlinespace
PEC observation (2023 Accreditation Visit): Lack of correlation in the survey forms used to assess certain PEO-2/PEO-3 sub-KPIs, reinforced by the IAB's concern over Alumni Survey Q5 (see above). & PEC Accreditation Visit, 2023, reiterated IAB Meeting, 13th November 2023 & The Alumni and Employer Survey instruments were fully redesigned, with items mapped directly to PEOs, as the Revised (2026) Evaluation System. The revised forms were reviewed and approved by the SE Departmental Board of Studies (DBS) on 20th April 2026. These were presented to and ratified by the SEECS HoDs' forum and the Faculty Board of Studies (FBS)of 21st April 2026, and finally presented to and endorsed by the IAB on 5th May 2026. The approved forms were then circulated by email to Alumni and Employers for input and used at the department's Open House in May 2026 (Appendix~\ref{app:governance_minutes}, Appendix~\ref{app:iab_minutes}).\\
\addlinespace
Introduce a Quantum Computing course. & IAB Meeting, 28th February 2025 (CEO, Emumba) & A dedicated elective, ``Introduction to Quantum Computing'' (offered in 3+1 and 2+1 credit-hour variants, covering superposition, entanglement, quantum algorithms, and hands-on labs using Qiskit-based simulators), was proposed via a SEECS-wide working paper (sponsored by SEECS, not the SE department alone) for inclusion as an elective in the BESE (Software Engineering), BSCS, BSAI, and BSDS curricula. The working paper records the course as approved by the DBS, for introduction from Spring 2026 (Appendix~\ref{app:governance_minutes}).\\
\addlinespace
Enable cross-departmental minor courses (e.g., SEECS students taking courses at NBS). & IAB Meeting, 28th February 2025 & Already available under NUST's University-wide Minors policy: SEECS students including (Software Engineering) may register for the Management minor at NUST Business School (NBS), subject to a minimum CGPA of 2.75, while SEECS reciprocally offers minors in Software Engineering, Computer Science, Data Science, and Artificial Intelligence to students of other NUST institutions (Appendix~\ref{app:governance_minutes}) (NUST Undergraduate Student Handbook, 28th August 2024 edition, Chapter 6, Para 2 ``Option of Minors'').\footnote{National University of Sciences and Technology, \textit{Undergraduate Student Handbook}, 28th August 2024 edition, Chapter 6, Para 2. Available at \url{https://seecs.nust.edu.pk/wp-content/uploads/2024/09/28-Aug-2024-Undergraduate-Handbook.pdf}.}\\
\addlinespace
Consider open-book examination formats. & IAB Meeting, 28th February 2025 (ITnIS Consulting) & Already available under existing NUST examination policy: the concerned faculty/invigilator may announce, in advance of a quiz, OHT, or mid/end-semester examination, the books, notes, or other material students may refer to during the test, effectively enabling an open-book format at faculty discretion (NUST Undergraduate Student Handbook, 28th August 2024 edition, Chapter 2, Para 8 ``Use of Reference Material during Test/Examination'').\footnote{National University of Sciences and Technology, \textit{Undergraduate Student Handbook}, 28th August 2024 edition, Chapter 2, Para 8. Available at \url{https://seecs.nust.edu.pk/wp-content/uploads/2024/09/28-Aug-2024-Undergraduate-Handbook.pdf}.}\\
\addlinespace
Introduce certification-style courses. & IAB Meeting, 28th February 2025 (ITnIS Consulting) & \textcolor{red}{NUST recently approved 1–3 CH range (per WP No. 2, 73rd ACM, 9th February 2026, see Appendix~\ref{app:certifications}) in consistence with HEC's CS/IT-specific elective-replacement provision (Letter HEC/CURR/GEN/2025/8036, 15th September 2025, Annex-A, Para 5), but differs from PEC's own proposal of a broader 6–12 CH range floated at its Task Force's inaugural meeting of 21st May 2025 and subsequently circulated for institutional input via Letter PEC/EAD/EQEEP/2025/04 (14th July 2025). That proposal remains subject to ratification by PEC's Engineering Accreditation Board (EAB) and Governing Body and has not yet been issued as a final SOP. The department intends to therefore adopt NUST policy in line with the more specific and more recent HEC guidance while awaiting PEC's final mechanism, and will revise credit-hour equivalencies if required once PEC's SOP is issued.}\\
\end{longtable}
\endgroup

\section{Standing and Emerging Recommendations (Full Record)}

\begingroup
\footnotesize
\setlength{\tabcolsep}{4pt}
\renewcommand{\arraystretch}{1.2}
\begin{longtable}[]{@{}>{\raggedright\arraybackslash}p{0.68\textwidth}>{\raggedright\arraybackslash}p{0.24\textwidth}@{}}
\caption{Open IAB Recommendations Pending Implementation}\label{tab:iab_open}\\
\toprule
\textbf{Recommendation (Not Yet Closed)} & \textbf{Source / Date}\\
\midrule
\endfirsthead
\toprule
\textbf{Recommendation (Not Yet Closed)} & \textbf{Source / Date}\\
\midrule
\endhead
\midrule
\multicolumn{2}{r}{\footnotesize\itshape Continued on next page}\\
\endfoot
\bottomrule
\endlastfoot
Expand coverage of AI agents, reverse engineering, and data-science algorithms. & IAB Meeting, 28th February 2025 (CEO, Emumba)\\
Establish a faculty ``boot camp'' to close the skills gap between academia and industry practice, targeted at PhD/MS scholars. & IAB Meeting, 7th May 2025, recorded as Action Point 2\\
\textcolor{red}{Broaden ``Applied AI in Education'' into a cross-sector applied-AI offering. Split DevOps and Machine Learning into separate courses.} & IAB Meeting, 7th May 2025, recorded as Action Point 1\\
Formalise a mentorship partnership (e.g., with Novatec X) to provide affordable, industry-led mentorship to students and researchers. & IAB Meeting, 7th May 2025, recorded as Action Point 3\\
Improve student and faculty access to premium cloud/AI tools to maximise use of available cloud-credit programmes. & IAB Meeting, 7th May 2025, recorded as Action Point 4\\
\textcolor{red}{Increase the number of internship placements and recruitment drives connected to Open House events.} & IAB Meeting, 28th February 2025\\
Continue to strengthen graduate soft skills, communication, customer orientation, and professional ethics, a recurring, sustained industry request across multiple sessions. & IAB Meetings: 24th May 2023, 2nd May 2024, 28th February 2025, 5th May 2026\\
\end{longtable}
\endgroup
